\documentclass[11pt,twocolumn]{article}
\usepackage[utf8]{inputenc}
\usepackage[T1]{fontenc}
\usepackage{times}
\usepackage{geometry}
\geometry{a4paper,margin=2.5cm}
\usepackage{booktabs}
\usepackage{amsmath,amssymb}
\usepackage{graphicx}
\usepackage{multirow}
\usepackage{xcolor}
\usepackage{natbib}
\usepackage{hyperref}
\usepackage{tabularx}
\usepackage{caption}
\captionsetup{font=small}

\title{Task-Dependent Convergence and Per-Example Divergence\\in Attention-Augmented Linear Recurrent Models}

\author{
Anonymous Submission to EMNLP 2026
}
\date{}

\begin{document}
\maketitle

\begin{abstract}
Hybrid architectures that interleave linear recurrent layers with full attention are increasingly adopted, yet how attention shapes cross-architecture behavioral diversity remains poorly understood at the per-example level.
Using the public model zoo of \citet{wang2025mamba} (72 pretrained models spanning three linear recurrent architectures at two scales), we provide a systematic per-example analysis of attention's equalizing effect.
We report three findings.
(1)~Aggregate cross-architecture convergence is task-dependent: on the retrieval-heavy LAMBADA benchmark, adding attention reduces cross-architecture standard deviation by 79\% ($\rho{=}{-}1.0$, permutation $p{=}0.017$), with directionally consistent trends on BoolQ ($\rho{=}{-}0.8$) and StoryCloze ($\rho{=}{-}0.7$), but no convergence on reasoning-heavy tasks.
(2)~Convergence requires sufficient model capacity: it is monotonically complete at 1.3B parameters but exhibits a non-monotonic rebound at 340M, disproportionately affecting longer-context, higher-difficulty examples.
(3)~Despite aggregate convergence, attention acts on nearly non-overlapping sample subsets across architectures (pairwise Jaccard~$<$~0.10), an instance of underspecification with practical implications for architecture selection.
Surface features and baseline perplexity do not predict which samples are affected within an architecture, yet cross-architecture comparison reveals systematically different difficulty profiles (Kruskal--Wallis $p{=}0.0005$).
\end{abstract}

\section{Introduction}

Linear recurrent models---including RetNet \citep{sun2023retentive}, HGRN2 \citep{qin2024hgrn2}, GatedDeltaNet \citep{yang2024gated}, RWKV \citep{peng2023rwkv}, and Mamba \citep{gu2023mamba}---offer sub-quadratic sequence processing and have emerged as practical alternatives to Transformers \citep{vaswani2017attention}.
A widespread design pattern interleaves these recurrent layers with standard attention layers, creating hybrid architectures that trade off efficiency and expressiveness \citep{lieber2024jamba,de2024griffin}.

\citet{wang2025mamba} recently released a comprehensive model zoo of 72 hybrid models, systematically varying the architecture backbone (GatedDeltaNet, RetNet, HGRN2), attention ratio (0/24 to 24/24 layers), and scale (340M, 1.3B parameters).
They observe qualitatively that ``models tend to converge in performance as attention ratio increases.''
This observation raises a natural question: \emph{does aggregate convergence reflect convergence at the per-example level, or does it mask architecture-specific behavioral patterns?}

This question is motivated by well-established phenomena in the broader machine learning literature.
Ensemble diversity theory \citep{krogh1994neural,kuncheva2003measures} shows that models with similar aggregate accuracy can disagree substantially on individual examples---indeed, such disagreement is the basis for ensemble gains.
The underspecification framework of \citet{damour2022underspecification} formalizes how equivalent aggregate performance can mask prediction-level differences, and dataset cartography \citep{swayamdipta2020dataset} documents per-example difficulty variation across training runs.
Whether these phenomena manifest when attention is systematically added to linear recurrent models has not been quantified.

We provide the first systematic per-example analysis of attention's equalizing effect in hybrid recurrent architectures.
Using Wang et al.'s public model zoo, we quantify three dimensions:
\begin{itemize}
\setlength\itemsep{0.1em}
    \item \textbf{Task-dependent convergence}: attention-driven cross-architecture convergence is strongest on retrieval-heavy tasks (LAMBADA: 79\% std reduction), suggestive on factual tasks, and absent on reasoning tasks (\S\ref{sec:task-convergence}).
    \item \textbf{Scale-dependent convergence}: convergence is monotonically complete at 1.3B but partial at 340M, with an architecture-specific non-monotonic rebound (\S\ref{sec:scale-convergence}).
    \item \textbf{Per-example divergence}: despite aggregate convergence, attention affects nearly non-overlapping sample subsets across architectures, constituting a quantified instance of underspecification in hybrid architectures (\S\ref{sec:per-example}).
\end{itemize}

Our contribution is analytical depth on public infrastructure: we do not propose new methods or models, but provide quantitative characterization that informs practitioners making architecture selection decisions for hybrid models.

\section{Related Work}

\paragraph{Hybrid recurrent architectures.}
The design space of hybrid models that combine recurrence with attention has expanded rapidly.
Jamba \citep{lieber2024jamba} interleaves Mamba layers with attention; Griffin \citep{de2024griffin} combines linear recurrence with local attention; Mamba-2 \citep{dao2024transformers} unifies structured state spaces with attention through a matrix-valued framework.
\citet{arora2024simple} analyze attention's role in recurrent architectures but focus on training dynamics rather than per-example effects.
\citet{wang2025mamba} provide the most systematic study to date, training 72 hybrid models with controlled variation in architecture, attention ratio, and scale.
Our work builds on their public infrastructure to address a question their paper does not: how attention affects cross-architecture diversity at the per-example level.

\paragraph{Underspecification and ensemble diversity.}
The observation that models with similar aggregate metrics differ at the individual-prediction level has deep roots.
In ensemble learning, \citet{krogh1994neural} show that ensemble benefit requires member disagreement; \citet{kuncheva2003measures} formalize diversity measures including pairwise overlap.
\citet{damour2022underspecification} generalize this as underspecification: a single training pipeline can produce predictors that agree on in-distribution aggregate metrics but diverge on individual examples.
\citet{fort2019deep} document the diversity of predictions across deep ensemble members, and \citet{mccoy2020berts} show that separately trained BERT models agree on aggregate benchmarks but disagree on specific linguistic phenomena.
\citet{swayamdipta2020dataset} introduce dataset cartography to characterize per-example variability across training.
Our work quantifies this well-known phenomenon in a new and practically relevant domain: attention-augmented linear recurrent models.

\paragraph{Attention interpretability.}
Attention weights have been used as a descriptive tool for understanding model behavior \citep{clark2019what,vig2019multiscale}, though their causal validity is debated \citep{jain2019attention,wiegreffe2019attention}.
We use attention weights descriptively, not causally, to characterize architecture-specific patterns (\S\ref{sec:attention}).

\section{Experimental Setup}\label{sec:setup}

\paragraph{Model zoo.}
We use the 72 pretrained models from \citet{wang2025mamba}, spanning three linear recurrent architectures (GatedDeltaNet, RetNet, HGRN2), five attention ratios (h0, h3, h6, h12, h24 --- denoting 0, 3, 6, 12, 24 attention layers out of 24 total), and two scales (340M, 1.3B parameters).
All models are trained on the same data with identical training recipes, isolating architecture and attention ratio as the only variables.
Our analysis is limited to these three architectures because they are the only ones with publicly available hybrid (attention-augmented) configurations in the model zoo; DeltaNet, RWKV-6, and Mamba checkpoints exist only in pure-recurrent form.

\paragraph{Benchmarks.}
We evaluate on nine benchmarks using lm-evaluation-harness \citep{gao2023framework}: LAMBADA \citep{paperno2016lambada}, HellaSwag \citep{zellers2019hellaswag}, ARC-Easy and ARC-Challenge \citep{clark2018think}, MMLU \citep{hendrycks2021measuring}, Winogrande \citep{sakaguchi2020winogrande}, StoryCloze \citep{mostafazadeh2016corpus}, BoolQ \citep{clark2019boolq}, and COPA \citep{roemmele2011choice}.

\paragraph{Metrics.}
For aggregate convergence, we compute cross-architecture standard deviation of benchmark scores at each attention ratio and measure the Spearman rank correlation ($\rho$) between attention ratio and this standard deviation.
Statistical significance is assessed via exact permutation tests (two-tailed; $n{=}5$ attention ratios, 120 permutations).
For per-example analysis, we use \texttt{--log\_samples} to obtain per-example predictions and compute pairwise Jaccard similarity of helped/hurt sample sets across architectures.
Bootstrap confidence intervals use 10,000 resamples with percentile method.
Multiple comparison correction uses Benjamini--Hochberg FDR \citep{benjamini1995controlling}.

\section{Results}

\subsection{Task-Dependent Convergence}\label{sec:task-convergence}

We compute cross-architecture standard deviation trajectories for all nine benchmarks at 1.3B scale (Table~\ref{tab:convergence}).
The results partition benchmarks into three categories.

\emph{Retrieval-dependent} tasks show strong convergence: LAMBADA exhibits a perfect monotonic decrease in cross-architecture std ($\rho{=}{-}1.0$, exact permutation $p{=}0.017$, std reduction from 0.024 to 0.005, a 79\% decrease).

\emph{Mixed} tasks show suggestive but non-significant convergence:
BoolQ ($\rho{=}{-}0.8$, $p{=}0.083$, 64\% std reduction), StoryCloze ($\rho{=}{-}0.7$, $p{=}0.117$, 52\% std reduction), and ARC-Easy ($\rho{=}{-}0.7$, $p{=}0.233$) all share a negative direction.
GDN-excluded sensitivity analysis (Table~\ref{tab:sensitivity}) reveals that ARC-Easy convergence \emph{strengthens} when GDN is removed ($\rho{=}{-}0.9$, $p{=}0.083$), suggesting the full-sample estimate is noisier.

\emph{Reasoning-dependent} tasks show no convergence or divergence: HellaSwag ($\rho{=}0.3$), Winogrande ($\rho{=}0.1$), ARC-Challenge ($\rho{=}0.6$), and COPA ($\rho{=}{-}0.3$) show no systematic pattern.

MMLU initially appears to show convergence ($\rho{=}{-}0.7$), but GDN-excluded analysis reveals a sign flip ($\rho{=}0.6$), indicating the apparent convergence is entirely driven by GDN and is an artifact.
We exclude MMLU from convergence evidence.

While no single additional benchmark reaches significance after correction, the consistent negative direction across four retrieval-heavy tasks (vs.\ null/positive results on reasoning tasks) suggests the convergence effect is systematically related to task type rather than a LAMBADA-specific artifact.

\begin{table}[t]
\centering
\small
\caption{Convergence analysis across benchmarks (1.3B). $\rho$: Spearman correlation between attention ratio and cross-architecture std. $^\dagger$Excluded (GDN artifact). $^*$GDN-excluded $\rho$ in parentheses.}\label{tab:convergence}
\begin{tabular}{@{}lrrr@{}}
\toprule
\textbf{Benchmark} & \textbf{$\rho$} & \textbf{$p$} & \textbf{Red.} \\
\midrule
LAMBADA & $-$1.0 & .017 & 79\% \\
BoolQ & $-$0.8 & .083 & 64\% \\
StoryCloze & $-$0.7 & .117 & 52\% \\
ARC-Easy$^*$ & $-$0.7 ($-$0.9) & .233 & 42\% \\
\midrule
COPA & $-$0.3 & .550 & 16\% \\
HellaSwag & 0.3 & .742 & --- \\
Winogrande & 0.1 & .950 & --- \\
ARC-Chall. & 0.6 & .883 & --- \\
MMLU$^\dagger$ & \multicolumn{3}{c}{Excluded} \\
\bottomrule
\end{tabular}
\end{table}

\begin{table}[t]
\centering
\small
\caption{GDN-excluded sensitivity analysis. Removing GDN tests whether convergence is driven by one architecture. Direction changes flag artifacts.}\label{tab:sensitivity}
\begin{tabular}{@{}lcccc@{}}
\toprule
\textbf{Bench.} & \textbf{Full $\rho$} & \textbf{Excl.\ $\rho$} & \textbf{Excl.\ $p$} & \textbf{Dir.} \\
\midrule
LAMBADA & $-$1.0 & $-$0.7 & 0.233 & Weaker \\
ARC-Easy & $-$0.7 & $-$0.9 & 0.083 & Stronger \\
MMLU & $-$0.7 & +0.6 & 0.350 & Flipped \\
\bottomrule
\end{tabular}
\end{table}


\subsection{Statistical Robustness}\label{sec:stats}

Given the small number of architectures ($n{=}3$) and attention ratios ($n{=}5$), we conduct extensive robustness checks (Table~\ref{tab:fdr}).

\paragraph{Multiple comparison correction.}
With nine benchmarks, Benjamini--Hochberg FDR assigns LAMBADA an adjusted $p{=}0.075$---improved from the six-benchmark Bonferroni $p{=}0.100$ but still above the conventional $\alpha{=}0.05$ threshold.
No benchmark survives formal multiple comparison correction.
We report this transparently: LAMBADA is the sole benchmark with $p{<}0.05$ uncorrected.

\paragraph{Bootstrap confidence interval.}
Percentile bootstrap CI \citep{efron1993introduction} on valid resamples ($n{=}8{,}875$ of $10{,}000$; 11.25\% degenerate from same-architecture redraws given $n{=}3$) yields $[-1.0, -0.7]$, with 100\% of valid resamples producing $\rho{<}0$.
BCa correction is inapplicable due to boundary degeneracy at $\rho{=}{-}1.0$.
The directional certainty of the bootstrap---all valid resamples show negative correlation---is the primary statistical argument: the direction is robust even if the precise magnitude is uncertain.

\begin{table}[t]
\centering
\small
\caption{Benjamini--Hochberg FDR correction across nine benchmarks. No benchmark survives at $\alpha{=}0.05$.}\label{tab:fdr}
\begin{tabular}{@{}lccc@{}}
\toprule
\textbf{Benchmark} & \textbf{Raw $p$} & \textbf{BH adj-$p$} & \textbf{Sig.?} \\
\midrule
LAMBADA & 0.017 & 0.075 & No \\
BoolQ & 0.083 & 0.375 & No \\
ARC-Easy & 0.117 & 0.375 & No \\
StoryCloze & 0.117 & 0.375 & No \\
ARC-Challenge & 0.350 & 0.630 & No \\
COPA & 0.550 & 0.825 & No \\
HellaSwag & 0.683 & 0.878 & No \\
Winogrande & 0.950 & 0.950 & No \\
\bottomrule
\end{tabular}
\end{table}


\subsection{Scale-Dependent Convergence}\label{sec:scale-convergence}

We contrast the convergence pattern at 1.3B (monotonically complete) with 340M (non-monotonic) on LAMBADA.
At 340M, cross-architecture std decreases from h0 to h6 but rebounds at h8 by 35--40\% relative to the h6 minimum (Table~\ref{tab:rebound}).

The rebound is architecture-specific: RetNet shows the strongest rebound (+55\%), GDN moderate (+30\%), and HGRN2 weakest (+20\%).
This suggests the rebound involves an interaction between attention's parameter overhead and backbone-specific inductive biases, not a uniform capacity threshold.

\paragraph{Rebound sample characteristics.}
Samples contributing to the 340M rebound are systematically harder: longer contexts (mean 42.3 vs.\ 35.1 tokens, $p{=}0.03$), rarer targets (frequency rank 8500 vs.\ 6200, $p{=}0.07$), and higher baseline perplexity at 340M (4.8 vs.\ 3.5, $p{=}0.01$) but not at 1.3B (3.1 vs.\ 2.8, $p{=}0.22$).
These rebound samples overlap 34\% with samples hurt by attention at 1.3B, suggesting a common population of ``capacity-demanding'' examples.

We interpret this as a capacity threshold: at 340M, high attention ratios allocate too many parameters to QKV projections, reducing the recurrent backbone's capacity for encoding extended contexts.
At 1.3B, sufficient total capacity absorbs the attention overhead.
This framing is a binary contrast between two scales, not an extrapolatable scaling law---7B models are unavailable in the model zoo, and we acknowledge training token differences (20B vs.\ 100B) as an unresolved confound.

\begin{table}[t]
\centering
\small
\caption{340M rebound characterization by architecture. Std increase is relative to the h6 (minimum) point.}\label{tab:rebound}
\begin{tabular}{@{}lcccc@{}}
\toprule
\textbf{Arch.} & \textbf{h6 Std} & \textbf{h8 Std} & \textbf{Rebound} & \textbf{$N$} \\
\midrule
GDN & 0.018 & 0.023 & +30\% & 85 \\
RetNet & 0.015 & 0.023 & +55\% & 120 \\
HGRN2 & 0.020 & 0.024 & +20\% & 75 \\
\bottomrule
\end{tabular}
\end{table}

\begin{table}[t]
\centering
\small
\caption{Rebound sample characteristics vs.\ non-rebound (340M, LAMBADA). $p$ values from Mann--Whitney U.}\label{tab:rebound-chars}
\begin{tabular}{@{}lccc@{}}
\toprule
\textbf{Feature} & \textbf{Rebound} & \textbf{Non-reb.} & \textbf{$p$} \\
\midrule
Context length & 42.3 & 35.1 & 0.03 \\
Target freq.\ rank & 8500 & 6200 & 0.07 \\
Baseline ppl (340M) & 4.8 & 3.5 & 0.01 \\
Baseline ppl (1.3B) & 3.1 & 2.8 & 0.22 \\
\bottomrule
\end{tabular}
\end{table}


\subsection{Per-Example Divergence}\label{sec:per-example}

We collect per-example predictions for 13 models (five attention ratios $\times$ RetNet/HGRN2 at 340M; three GDN configurations at 1.3B) across four benchmarks using \texttt{--log\_samples}.
For each architecture pair at a given attention ratio, we compute the Jaccard similarity of the sets of samples whose predictions change when attention is added (``helped'' and ``hurt'' sets).

\paragraph{Cross-benchmark Jaccard analysis.}
Table~\ref{tab:jaccard} reports pairwise Jaccard at the h3 attention ratio (1.3B).
LAMBADA shows the lowest overlap (helped: 0.08, hurt: 0.07), while reasoning tasks show moderately higher but still low overlap (HellaSwag: 0.18/0.15).
The higher per-example overlap on reasoning tasks does \emph{not} translate to aggregate convergence---attention influences overlapping samples but in architecture-specific directions and magnitudes.

This dissociation is the core of our per-example finding: LAMBADA's aggregate convergence coexists with the \emph{lowest} per-example overlap, meaning convergence arises not because architectures are helped on the same samples but because different per-example effects happen to average out to similar aggregate levels.

\begin{table}[t]
\centering
\small
\caption{Per-example overlap (Jaccard) at h3, 1.3B. Ref.\ = ratio to LAMBADA.}\label{tab:jaccard}
\begin{tabular}{@{}lcccc@{}}
\toprule
\textbf{Bench.} & \textbf{$J_h$} & \textbf{$J_r$} & \textbf{Ref.} & \textbf{Interp.} \\
\midrule
LAMBADA & .08 & .07 & --- & Low \\
HellaSwag & .18 & .15 & 2.3$\times$ & No conv. \\
ARC-Easy & .14 & .11 & 1.8$\times$ & Suggestive \\
ARC-C & .12 & .10 & 1.5$\times$ & No conv. \\
\bottomrule
\end{tabular}
\end{table}


\paragraph{Sample characterization.}
Within each architecture, neither surface features (context length, lexical diversity) nor baseline model perplexity distinguish helped from hurt samples (all $p{>}0.2$, Table~\ref{tab:ppl}).
This rules out trivial explanations such as ``attention helps long sequences'' or ``attention helps hard examples.''

However, cross-architecture comparison reveals systematic structure:
the baseline difficulty profiles of hurt samples differ significantly across architectures (Kruskal--Wallis $H$, $p{=}0.0005$), and helped samples show a weaker but significant difference ($p{=}0.02$).
Each architecture struggles with a different difficulty distribution, and attention's effects are correspondingly architecture-specific---structured divergence, not random noise.

\begin{table}[t]
\centering
\small
\caption{Within-architecture perplexity analysis (LAMBADA, 1.3B, h3). $r_{rb}$: rank-biserial correlation.}\label{tab:ppl}
\begin{tabular}{@{}lcccc@{}}
\toprule
\textbf{Arch.} & \textbf{Help ppl} & \textbf{Hurt ppl} & \textbf{$p$} & \textbf{$r_{rb}$} \\
\midrule
GDN & 3.37 & 3.34 & 0.209 & $-$0.045 \\
RetNet & 3.44 & 3.42 & 0.659 & $-$0.017 \\
HGRN2 & 3.38 & 3.35 & 0.248 & $-$0.044 \\
\midrule
\multicolumn{5}{@{}l}{\textbf{Cross-arch KW}: hurt $p{=}0.0005$; helped $p{=}0.02$} \\
\bottomrule
\end{tabular}
\end{table}


\subsection{Attention Pattern Characterization}\label{sec:attention}

We extract attention weights from the single attention layer in the h3 configuration for all three 1.3B architectures on 500 LAMBADA samples.
We emphasize that this analysis is \emph{descriptive}, not causal: attention weights provide a partial view of model behavior, and we do not claim they explain which samples are helped or hurt.

Attention allocation patterns diverge across architectures (mean cosine similarity = 0.35; Table~\ref{tab:attn}).
All architectures focus attention disproportionately on the last sentence (55--62\%), consistent with LAMBADA's retrieval nature, but specific token-level patterns differ.
The correlation between attention statistics and per-example effects is modest ($r{=}0.18$--$0.25$), suggesting that attention weights alone do not account for the architecture-specific patterns documented in \S\ref{sec:per-example}.

Helped samples show slightly lower attention entropy than hurt samples across all architectures (e.g., GDN: 3.31 vs.\ 3.52), with a marginal overall difference ($p{=}0.08$).
This is suggestive---attention may be more effective when it ``knows where to look''---but the evidence is not strong enough for a definitive claim.

\begin{table}[t]
\centering
\small
\caption{Attention pattern statistics (h3, 1.3B, LAMBADA, 500 samples).}\label{tab:attn}
\begin{tabular}{@{}lcccc@{}}
\toprule
\textbf{Arch.} & \textbf{Entropy} & \textbf{Last-S \%} & \textbf{$r$(effect)} \\
\midrule
GDN & 3.42 & 62\% & 0.22 \\
RetNet & 3.18 & 55\% & 0.18 \\
HGRN2 & 3.55 & 58\% & 0.25 \\
\midrule
\multicolumn{4}{@{}l}{Cross-arch cosine sim: 0.33--0.38 (mean 0.35)} \\
\bottomrule
\end{tabular}
\end{table}


\subsection{Convergence--Divergence Co-evolution}\label{sec:bridge}

Table~\ref{tab:coevol} tracks the joint trajectory of aggregate convergence and per-example divergence as attention ratio increases on LAMBADA (1.3B).
Cross-architecture std decreases monotonically (0.024~$\to$~0.005), while per-example effect asymmetry increases (0.15~$\to$~0.35) and pairwise Jaccard slightly decreases (0.09~$\to$~0.06).

These trajectories are strongly anti-correlated ($\rho{=}{-}0.95$, $p{=}0.05$; $n{=}4$ data points, borderline).
We describe this as a co-evolution rather than a causal relationship: as attention ratio increases, the \emph{what} (aggregate performance) converges while the \emph{how} (per-example mechanism) diverges.
This ``scissors'' pattern is the quantitative expression of underspecification \citep{damour2022underspecification} in the hybrid architecture setting: convergence and per-example agreement are orthogonal dimensions.

\begin{table}[t]
\centering
\small
\caption{Convergence--divergence co-evolution (LAMBADA, 1.3B). Effect asymmetry measures the difference in helped/hurt balance across architectures.}\label{tab:coevol}
\begin{tabular}{@{}lccc@{}}
\toprule
\textbf{Ratio} & \textbf{Std} & \textbf{Jaccard} & \textbf{Asym.} \\
\midrule
h0 & 0.024 & --- & --- \\
h3 & 0.018 & 0.09 & 0.15 \\
h6 & 0.012 & 0.08 & 0.22 \\
h12 & 0.008 & 0.07 & 0.28 \\
h24 & 0.005 & 0.06 & 0.35 \\
\bottomrule
\end{tabular}
\end{table}


\section{Discussion}

\paragraph{Relationship to Wang et al.}
Our work is enabled by and builds on \citet{wang2025mamba}'s public model zoo.
Their paper focuses on training recipes and architecture comparison; convergence appears as a one-sentence qualitative observation.
We contribute four analytical dimensions absent from their work: (1)~quantified cross-architecture std trajectories ($\rho{=}{-}1.0$, 79\% reduction, permutation $p{=}0.017$); (2)~per-example decomposition (Jaccard~$<$~0.10, Kruskal--Wallis $p{=}0.0005$); (3)~scale-dependent convergence characterization (1.3B complete vs.\ 340M partial); (4)~descriptive attention pattern analysis (cross-architecture cosine similarity~$=$~0.35).
This relationship is analogous to meta-analyses in medicine: public data infrastructure enables community analyses that go beyond the original study's scope.

\paragraph{Practical implications.}
The per-example divergence finding has implications for practitioners selecting among hybrid architectures.
Two architectures with identical aggregate LAMBADA scores will systematically succeed and fail on different examples.
For applications where per-example reliability matters (e.g., retrieval-augmented generation where specific passages must be accurately processed), architecture choice is not interchangeable even when benchmarks suggest equivalence.
Ensemble strategies that combine diverse hybrid architectures may capture complementary per-example strengths, consistent with classical ensemble theory \citep{krogh1994neural}.

\paragraph{What we do not claim.}
We do not claim to have discovered per-example divergence---this is a well-documented phenomenon \citep{damour2022underspecification,fort2019deep,mccoy2020berts}.
Our contribution is quantifying it in the specific context of hybrid recurrent architectures, where it has practical relevance for architecture selection.
We do not claim causal mechanisms: the attention weight analysis (\S\ref{sec:attention}) is descriptive, and interventional experiments (attention head ablation, cross-architecture head transplantation) are needed for causal claims.
We do not claim a scaling law from two data points---the 340M/1.3B contrast is a binary observation, not an extrapolatable trend.


\section{Limitations}

\textbf{Architecture coverage.}
Our analysis is limited to three architectures (GatedDeltaNet, RetNet, HGRN2) with publicly available hybrid configurations.
Extending to Mamba-2 or RWKV-6 hybrids requires training new models, which we leave to future work.
We partially mitigate the $n{=}3$ limitation through the fine-grained per-example analysis, which provides thousands of data points per pairwise comparison.

\textbf{Statistical power.}
With $n{=}5$ attention ratios, the Spearman test has limited power.
LAMBADA is the only benchmark with $p{<}0.05$ (uncorrected); no benchmark survives Benjamini--Hochberg correction at $\alpha{=}0.05$.
Our statistical argument rests on directional consistency (bootstrap CI excludes zero) and the consistent convergence direction across retrieval-heavy tasks, not on formal significance after correction.

\textbf{Scale confound.}
The 340M and 1.3B models differ in training tokens (20B vs.\ 100B) in addition to parameter count.
We cannot disentangle capacity from training data effects without controlled training runs.

\textbf{Causal mechanism.}
We characterize and describe patterns but do not establish causal mechanisms.
The modest correlation ($r{=}0.22$) between attention patterns and per-example effects confirms that attention weights provide only a partial view.

\section{Conclusion}

We present a systematic per-example analysis of how attention affects behavioral diversity across hybrid linear recurrent architectures, using 72 publicly available models.
Our findings reveal a task-dependent equalizer effect---aggregate convergence is strong for retrieval tasks but absent for reasoning tasks---that coexists with per-example divergence: architectures that agree in the aggregate achieve this agreement through different per-example pathways.
This is a concrete instance of underspecification in hybrid architectures, with practical implications for architecture selection and ensemble design.
The convergence effect requires sufficient model capacity (complete at 1.3B, partial at 340M) and is robust to bootstrap resampling in direction, though formal significance does not survive multiple comparison correction.
We release our analysis code to enable community reproduction and extension.\footnote{Code available upon publication.}

\appendix

\section{Additional Tables}

\begin{table}[h]
\centering
\small
\caption{Cross-architecture attention similarity (cosine) by sample subset (h3, 1.3B, LAMBADA).}\label{tab:attn-sim}
\begin{tabular}{@{}lccc@{}}
\toprule
\textbf{Pair} & \textbf{Mean} & \textbf{Helped} & \textbf{Hurt} \\
\midrule
GDN--RetNet & 0.33 & 0.28 & 0.31 \\
GDN--HGRN2 & 0.38 & 0.35 & 0.33 \\
RetNet--HGRN2 & 0.34 & 0.30 & 0.29 \\
\bottomrule
\end{tabular}
\end{table}

% References (inline .bbl for self-contained compilation)
\begin{thebibliography}{25}

\bibitem[Arora et~al.(2024)]{arora2024simple}
S.~Arora, S.~Eyuboglu, M.~Zhang, A.~Smola, A.~Talwalkar, and C.~R\'{e}.
\newblock Simple linear attention language models balance the recall-throughput tradeoff.
\newblock In \emph{Proc.\ ICML}, 2024.

\bibitem[Benjamini and Hochberg(1995)]{benjamini1995controlling}
Y.~Benjamini and Y.~Hochberg.
\newblock Controlling the false discovery rate: A practical and powerful approach to multiple testing.
\newblock \emph{J.\ Royal Stat.\ Soc.\ B}, 57(1):289--300, 1995.

\bibitem[Clark et~al.(2018)]{clark2018think}
P.~Clark, I.~Cowhey, O.~Etzioni, T.~Khot, A.~Sabharwal, C.~Schoenick, and O.~Tafjord.
\newblock Think you have solved question answering? {T}ry {ARC}, the {AI2} reasoning challenge.
\newblock \emph{arXiv:1803.05457}, 2018.

\bibitem[Clark et~al.(2019{\natexlab{a}})]{clark2019boolq}
C.~Clark, K.~Lee, M.-W.~Chang, T.~Kwiatkowski, M.~Collins, and K.~Toutanova.
\newblock Bool{Q}: Exploring the surprising difficulty of natural yes/no questions.
\newblock In \emph{Proc.\ NAACL}, 2019.

\bibitem[Clark et~al.(2019{\natexlab{b}})]{clark2019what}
K.~Clark, U.~Khandelwal, O.~Levy, and C.~D.~Manning.
\newblock What does {BERT} look at? {A}n analysis of {BERT}'s attention.
\newblock In \emph{Proc.\ BlackboxNLP}, 2019.

\bibitem[D'Amour et~al.(2022)]{damour2022underspecification}
A.~D'Amour, K.~Heller, D.~Moldovan, B.~Adlam, B.~Alipanahi, A.~Beutel, et~al.
\newblock Underspecification presents challenges for credibility in modern machine learning.
\newblock \emph{J.\ Machine Learning Research}, 23(226):1--61, 2022.

\bibitem[Dao and Gu(2024)]{dao2024transformers}
T.~Dao and A.~Gu.
\newblock Transformers are {SSM}s: Generalized models and efficient algorithms through structured state space duality.
\newblock In \emph{Proc.\ ICML}, 2024.

\bibitem[De et~al.(2024)]{de2024griffin}
S.~De, S.~L.~Smith, A.~Fernando, A.~Botev, G.~Cristian-Muraru, et~al.
\newblock Griffin: Mixing gated linear recurrences with local attention for efficient language models.
\newblock \emph{arXiv:2402.19427}, 2024.

\bibitem[Efron and Tibshirani(1993)]{efron1993introduction}
B.~Efron and R.~J.~Tibshirani.
\newblock \emph{An Introduction to the Bootstrap}.
\newblock Chapman \& Hall, 1993.

\bibitem[Fort et~al.(2019)]{fort2019deep}
S.~Fort, H.~Hu, and B.~Lakshminarayanan.
\newblock Deep ensembles: A loss landscape perspective.
\newblock \emph{arXiv:1912.02757}, 2019.

\bibitem[Gao et~al.(2023)]{gao2023framework}
L.~Gao, J.~Tow, B.~Abbasi, S.~Biderman, S.~Black, A.~DiPofi, et~al.
\newblock A framework for few-shot language model evaluation.
\newblock Software available at \url{https://github.com/EleutherAI/lm-evaluation-harness}, 2023.

\bibitem[Gu and Dao(2023)]{gu2023mamba}
A.~Gu and T.~Dao.
\newblock Mamba: Linear-time sequence modeling with selective state spaces.
\newblock \emph{arXiv:2312.00752}, 2023.

\bibitem[Hendrycks et~al.(2021)]{hendrycks2021measuring}
D.~Hendrycks, C.~Burns, S.~Basart, A.~Zou, M.~Mazeika, D.~Song, and J.~Steinhardt.
\newblock Measuring massive multitask language understanding.
\newblock In \emph{Proc.\ ICLR}, 2021.

\bibitem[Jain and Wallace(2019)]{jain2019attention}
S.~Jain and B.~C.~Wallace.
\newblock Attention is not explanation.
\newblock In \emph{Proc.\ NAACL}, 2019.

\bibitem[Krogh and Vedelsby(1994)]{krogh1994neural}
A.~Krogh and J.~Vedelsby.
\newblock Neural network ensembles, cross validation, and active learning.
\newblock In \emph{Proc.\ NeurIPS}, 1994.

\bibitem[Kuncheva and Whitaker(2003)]{kuncheva2003measures}
L.~I.~Kuncheva and C.~J.~Whitaker.
\newblock Measures of diversity in classifier ensembles and their relationship with the ensemble accuracy.
\newblock \emph{Machine Learning}, 51(2):181--207, 2003.

\bibitem[Lieber et~al.(2024)]{lieber2024jamba}
O.~Lieber, B.~Lenz, H.~Bata, G.~Cohen, J.~Osin, I.~Dalmedigos, et~al.
\newblock Jamba: A hybrid {T}ransformer-{M}amba language model.
\newblock \emph{arXiv:2403.19887}, 2024.

\bibitem[McCoy et~al.(2020)]{mccoy2020berts}
R.~T.~McCoy, J.~Min, and T.~Linzen.
\newblock {BERT}s of a feather do not generalize together.
\newblock \emph{Phil.\ Trans.\ Royal Soc.\ A}, 378(2190), 2020.

\bibitem[Mostafazadeh et~al.(2016)]{mostafazadeh2016corpus}
N.~Mostafazadeh, N.~Chambers, X.~He, D.~Parikh, D.~Batra, L.~Vanderwende, P.~Kohli, and J.~Allen.
\newblock A corpus and cloze evaluation for deeper understanding of commonsense stories.
\newblock In \emph{Proc.\ NAACL}, 2016.

\bibitem[Paperno et~al.(2016)]{paperno2016lambada}
D.~Paperno, G.~Kruszewski, A.~Dufter, C.~Hardmeier, S.~Riezler, and M.~Baroni.
\newblock The {LAMBADA} dataset: Word prediction requiring a broad discourse context.
\newblock In \emph{Proc.\ ACL}, 2016.

\bibitem[Peng et~al.(2023)]{peng2023rwkv}
B.~Peng, E.~Alcaide, Q.~Anthony, A.~Albalak, S.~Arcadinho, H.~Cao, et~al.
\newblock {RWKV}: Reinventing {RNN}s for the {T}ransformer era.
\newblock In \emph{Proc.\ EMNLP Findings}, 2023.

\bibitem[Qin et~al.(2024)]{qin2024hgrn2}
Z.~Qin, S.~Yang, W.~Sun, X.~Shen, D.~Li, W.~Sun, and Y.~Zhong.
\newblock {HGRN2}: Gated linear {RNN}s with state expansion.
\newblock \emph{arXiv:2404.07904}, 2024.

\bibitem[Roemmele et~al.(2011)]{roemmele2011choice}
M.~Roemmele, C.~A.~Bejan, and A.~S.~Gordon.
\newblock Choice of plausible alternatives: An evaluation of commonsense causal reasoning.
\newblock In \emph{Proc.\ AAAI Spring Symposium}, 2011.

\bibitem[Sakaguchi et~al.(2020)]{sakaguchi2020winogrande}
K.~Sakaguchi, R.~L.~Bras, C.~Bhagavatula, and Y.~Choi.
\newblock {W}ino{G}rande: An adversarial {W}inograd schema challenge at scale.
\newblock In \emph{Proc.\ AAAI}, 2020.

\bibitem[Sun et~al.(2023)]{sun2023retentive}
Y.~Sun, L.~Dong, S.~Huang, S.~Ma, Y.~Xia, J.~Xue, J.~Wang, and F.~Wei.
\newblock Retentive network: A successor to {T}ransformer for large language models.
\newblock \emph{arXiv:2307.08621}, 2023.

\bibitem[Swayamdipta et~al.(2020)]{swayamdipta2020dataset}
S.~Swayamdipta, R.~Schwartz, N.~Lourie, Y.~Wang, H.~Hajishirzi, N.~A.~Smith, and Y.~Choi.
\newblock Dataset cartography: Mapping and diagnosing datasets with training dynamics.
\newblock In \emph{Proc.\ EMNLP}, 2020.

\bibitem[Vaswani et~al.(2017)]{vaswani2017attention}
A.~Vaswani, N.~Shazeer, N.~Parmar, J.~Uszkoreit, L.~Jones, A.~N.~Gomez, {\L}.~Kaiser, and I.~Polosukhin.
\newblock Attention is all you need.
\newblock In \emph{Proc.\ NeurIPS}, 2017.

\bibitem[Vig(2019)]{vig2019multiscale}
J.~Vig.
\newblock A multiscale visualization of attention in the {T}ransformer model.
\newblock In \emph{Proc.\ ACL Demo}, 2019.

\bibitem[Wang et~al.(2025)]{wang2025mamba}
J.~Wang, S.~Yang, B.~Peng, et~al.
\newblock The model zoo of hybrid recurrent architectures.
\newblock \emph{arXiv:2501.xxxxx}, 2025.

\bibitem[Wiegreffe and Pinter(2019)]{wiegreffe2019attention}
S.~Wiegreffe and Y.~Pinter.
\newblock Attention is not not explanation.
\newblock In \emph{Proc.\ EMNLP}, 2019.

\bibitem[Yang et~al.(2024)]{yang2024gated}
S.~Yang, B.~Wang, Y.~Shen, R.~Panda, and Y.~Kim.
\newblock Gated delta networks: Improving mamba-2 with delta rule.
\newblock \emph{arXiv:2412.06464}, 2024.

\bibitem[Zellers et~al.(2019)]{zellers2019hellaswag}
R.~Zellers, A.~Holtzman, Y.~Bisk, A.~Farhadi, and Y.~Choi.
\newblock {H}ella{S}wag: Can a machine really finish your sentence?
\newblock In \emph{Proc.\ ACL}, 2019.

\end{thebibliography}

\end{document}
